\chapter{Timing Between the Chips}
\label{ch:synchronisation}

The three chips run from the same clock, so their timings can be stated against each other. This chapter gives the
relationships that a program or a tool needs: what the CPU can do in a line and in a frame, how the sound lines up with
the picture, and how long it takes to reach each chip.

\section{What the CPU Can Do in a Frame}
\label{sec:cpu-budget}

An ordinary instruction takes six \texttt{CLKI} cycles. Table~\ref{tab:cpu-budget} gives the number of instructions
that fit in each part of the picture. An instruction that waits on the bus, or a coprocessor call, takes longer
(\archmanual{}, instruction timing appendix).

\begin{table}[H]
	\centering
	\begin{tabular}{lrr}
		\toprule
		                                  & \textbf{NTSC} & \textbf{PAL} \\
		\midrule
		Cycles in a line                  & 682           & 682          \\
		Instructions in a line            & 113.7         & 113.7        \\
		Cycles in a frame                 & 178,684       & 212,784      \\
		Instructions in a frame           & 29,781        & 35,464       \\
		Cycles in vertical blanking       & 25,916        & 60,016       \\
		Instructions in vertical blanking & 4,319         & 10,003       \\
		\bottomrule
	\end{tabular}
	\caption{Instructions That Fit in the Picture}
	\label{tab:cpu-budget}
\end{table}

\noindent Those figures assume no wait states. Every instruction makes two bus accesses to be fetched, and each access to a
memory waits (Table~\ref{tab:wait-states}), so an ordinary instruction takes 6 cycles plus two for each wait state. A load or
a store adds one access more. The counts a program really gets are lower:

\begin{table}[H]
	\centering
	\begin{tabular}{lrrrr}
		\toprule
		                                  & \multicolumn{2}{c}{\textbf{From working memory or ROM}} & \multicolumn{2}{c}{\textbf{From the cartridge}}                                \\
		                                  & \textbf{NTSC}                                           & \textbf{PAL}                                    & \textbf{NTSC} & \textbf{PAL} \\
		\midrule
		Cycles an ordinary instruction    & 8                                                       & 8                                               & 10            & 10           \\
		Instructions in a line            & 85.3                                                    & 85.3                                            & 68.2          & 68.2         \\
		Instructions in a frame           & 22,336                                                  & 26,598                                          & 17,868        & 21,278       \\
		Instructions in vertical blanking & 3,240                                                   & 7,502                                           & 2,592         & 6,002        \\
		\bottomrule
	\end{tabular}
	\caption{Instructions That Fit in the Picture, With Wait States}
	\label{tab:cpu-budget-waits}
\end{table}

\noindent Code that runs from the cartridge goes at four fifths of the speed of code in working memory, so a program that needs the
speed copies its inner loops into working memory.

\noindent A program has the whole of vertical blanking to change the display adapter's memory, and the interrupt that
starts it arrives at its beginning (\videomanual{}, timing chapter).

\section{Sound Against Picture}
\label{sec:sound-picture}

SIRCAS makes one output sample every 336 cycles. A line is 682 cycles, so a line holds a little over two samples. The
numbers over a frame are:

\begin{table}[H]
	\centering
	\begin{tabular}{lrr}
		\toprule
		                                    & \textbf{NTSC} & \textbf{PAL} \\
		\midrule
		Audio samples in a line             & 2.030         & 2.030        \\
		Audio samples in a frame            & 531.80        & 633.29       \\
		Envelope updates in a frame         & 8.31          & 9.90         \\
		Sample periods in vertical blanking & 77.1          & 178.6        \\
		\bottomrule
	\end{tabular}
	\caption{Sound and Picture}
	\label{tab:sound-picture}
\end{table}

A frame does not hold a whole number of samples, and the sequencers count samples, not frames, so a sequence is never
locked to the picture. A program that changes the instruments of the sound once a frame does so at a different point of
the sample period each time, and that is not audible. A program that must keep a sequence in step with the picture
cannot rely on the two staying together, as the next section shows.

\subsection{Keeping a Sequencer Tick Near a Frame}

The tick of the sequencers is \texttt{SEQ\_PERIOD} samples (\audiomanual{}, sequencer chapter). The nearest whole
number of samples to a frame gives a tick close to a frame, and it drifts:

\begin{table}[H]
	\centering
	\begin{tabular}{lrr}
		\toprule
		                                      & \textbf{NTSC} & \textbf{PAL}  \\
		\midrule
		Frame, in samples                     & 531.80        & 633.29        \\
		\texttt{SEQ\_PERIOD} nearest a frame  & 532           & 633           \\
		Tick, in cycles                       & 178,752       & 212,688       \\
		Tick against the frame                & 0.038\% long  & 0.045\% short \\
		Frames for the tick to slip one frame & about 2,628   & about 2,216   \\
		Time for the tick to slip one frame   & 43.7 s        & 44.3 s        \\
		\bottomrule
	\end{tabular}
	\caption{A Tick of One Frame}
	\label{tab:tick-frame}
\end{table}

\noindent A program that needs the sound to stay with the picture over longer than that counts frames and changes the
period, or starts the sequence again, from time to time.

\section{Reaching the Chips}
\label{sec:reaching}

Table~\ref{tab:answer-times} gave the length of a bus cycle to each chip. What happens after a write differs between
them:

\begin{itemize}
	\item \textbf{SIRCADA.} A register write completes in the cycle it is made, and a video memory access takes two. The display adapter samples its registers
	      228 master clocks, 114 cycles, after the start of a line, so a write made earlier in the line than that takes
	      effect on the same line (\videomanual{}, timing chapter).
	\item \textbf{SIRCAS.} A register write also completes in the cycle. When the change is heard depends on what was
	      written: a note started by the CPU sounds a few sample periods later, about a line and a half of the picture,
	      and some changes wait for the next envelope update, which is up to 2.0~ms, about 31 lines, away
	      (\audiomanual{}, playback and envelope chapters).
	\item \textbf{The audio memory.} A word written through the window is in the memory when the cycle ends, and a
	      voice that reads that word reads the new value from then on. The window takes two to five cycles a word
	      (Table~\ref{tab:answer-times}).
\end{itemize}

\section{Electrical Margins}
\label{sec:margins}

The timing diagrams of the manuals show logical timing: which cycle an event falls in. Each access must also fit in
real time. The CPU puts out the address and \activelow{BAS} at the start of a cycle and samples \activelow{BACK} and
the read data at the rising edge that ends it. What is left for the CPU's own output delay and set-up time is what
remains after the decoder, the glue gate and the memory. Table~\ref{tab:margins} counts those against the cycles of the
access, with the delays of the 74ALS138 (20~ns), a glue gate (15~ns) and about 5~ns of board.

\begin{table}[H]
	\centering
	\small
	\begin{tabular}{lrrrrr}
		\toprule
		\textbf{Access}              & \textbf{Time (ns)} & \textbf{Decoder and gate} & \textbf{Board} & \textbf{Memory} & \textbf{Left (ns)} \\
		\midrule
		Working memory, 70~ns        & 186.2              & 35                        & 5              & 70              & 76.2               \\
		System ROM, 70~ns            & 186.2              & 35                        & 5              & 70              & 76.2               \\
		Cartridge, 150~ns            & 279.3              & 35                        & 10             & 150             & 84.3               \\
		Audio memory, 100~ns         & 186.2              & ---                       & 5              & 100             & 81.2               \\
		Video memory, 35~ns          & 93.1               & ---                       & 5              & 35              & 53.1               \\
		SIRCADA or SIRCAS, one cycle & 93.1               & 20                        & 5              & ---             & 68.1               \\
		\bottomrule
	\end{tabular}
	\caption{Time Left in an Access}
	\label{tab:margins}
\end{table}

\noindent What is left must cover the output delay of the device that drives the address and the set-up time of the device that
samples the data: that of the CPU for the first three rows, of SIRCAS and the SRAM for the audio memory, and of SIRCADA for
the video memory. For the last row it must also cover the device's own delay from \activelow{CS} to valid data and to
\activelow{BACK}. A 200~ns cartridge ROM would leave only 34~ns, and a 100~ns working memory 46~ns, so the table of wait states asks for
faster parts, to keep a margin. The decoder is the largest delay in the path: a 74LS138 would take 20~ns more from every row
that has one.

\begin{itemize}
	\item \textbf{Write strobe.} A static memory is written while its write enable is low, and needs the data to stay valid
	      after the strobe ends. The strobe must not end when \activelow{BAS} rises, because the CPU stops driving the data
	      about when it does, and the decoder's delay would make the strobe end after the data has gone. It is made instead to
	      end at the falling edge of \texttt{CLKI} in the last cycle of the access, so that it ends half a cycle (46~ns) before the
	      data is released, and the data is held for that long. It begins when the select is valid, about 40~ns into the
	      first cycle, so a two-cycle write has a strobe of about 110~ns and a three-cycle write of about 200~ns, both longer
	      than the 70~ns parts need.
	\item \textbf{Read hold.} A memory holds its data for a few nanoseconds after its select or address changes. The CPU
	      samples at the edge, and the select changes after the edge by its own delay and the decoder's, so the data is held
	      for at least the 20~ns of the decoder.
	\item \textbf{One-cycle devices.} SIRCADA and SIRCAS answer in one cycle. After the decoder, 68~ns of the cycle remains
	      for the CPU's output delay, the device's delay from \activelow{CS} to valid data and to \activelow{BACK}, and the
	      CPU's set-up time. The devices' own delay should therefore not exceed 35~ns.
	\item \textbf{Video memory.} SIRCADA's delay from clock to address, the SRAM's 35~ns and SIRCADA's data set-up time must
	      sum to no more than 88~ns (93.1~ns less 5~ns of board), so the delay and the set-up together must not exceed 53~ns.
\end{itemize}

\noindent The electrical figures of the chips themselves are not yet specified, so these margins are the requirements that the
chips must meet. The chips are designed to the following targets (the CPU's are in Chapter~\ref{ch:electrical-characteristics} of the \hwmanual{}), which every row of Table~\ref{tab:margins} accommodates:
the CPU's address, \activelow{BAS}, BRW and data outputs are valid within 25~ns of the rising edge of \texttt{CLKI}, its data and
\activelow{BACK} inputs need 10~ns of set-up and no hold, and SIRCADA's and SIRCAS's outputs and set-up times are the same.
The CPU's 35~ns leaves 41~ns on the memory paths and 33~ns on the one-cycle devices.

\section{Loading the Audio Memory}

Filling the audio memory is the slowest thing a program does to the system. The window takes two to five cycles to
accept a word, so a transfer of $n$ words takes between $2n$ and $5n$ cycles for the writes alone, on top of whatever
the CPU takes to fetch the words. Table~\ref{tab:load-times} gives the writes for a few sizes.

\begin{table}[H]
	\centering
	\begin{tabular}{rrrr}
		\toprule
		\textbf{Words} & \textbf{Fastest (cycles)} & \textbf{Slowest (cycles)} & \textbf{Slowest, in NTSC lines} \\
		\midrule
		256            & 512                       & 1,280                     & 1.9                             \\
		4,096          & 8,192                     & 20,480                    & 30.0                            \\
		32,768         & 65,536                    & 163,840                   & 240.2                           \\
		\bottomrule
	\end{tabular}
	\caption{Writing the Audio Memory Through the Window}
	\label{tab:load-times}
\end{table}

\noindent The whole memory therefore takes at most 163,840 cycles, 15.3~ms, to write: less than one NTSC frame, but far longer than the 2.41~ms of vertical blanking. The sound plays while it is being written, so a program loads it before it enables
the output (\audiomanual{}, programming chapter).
